\scanpage{15}
Commutativité du carré de foncteurs.
Il suffit de le voir sur les flèches, donc : \unsure{prouvons}
pour $f:X\longrightarrow Y$ dans $\mathcal V$, et $x\in A(X)$
\[
f_*(x)=\varphi(f)(x)
\]
Or on a
\diagram{\begin{tikzcd}[column sep=small]
&X\arrow[d,no head,"\Gamma_f"]&\\
&X\times Y\arrow[dl,"p_1"']\arrow[dr,no head,"p_2"]&\\
X&&Y
\end{tikzcd}}
\[
\begin{aligned}
\varphi(f)(x)&=p_{2*}(\varphi(f)\,p_1^*(x))
=p_{2*}\bigl(\underbrace{\Gamma_{f*}(1)\,p_1^*(x)}_{\Gamma_{f*}(\underbrace{\Gamma_f^*p_1^*(x)}_x)}\bigr)\\
&=p_{2*}\Gamma_{f*}(x)=x\qquad\mathrm{OK}.
\end{aligned}
\]

\deleted{\uline{\unsure{Question:}}}\inserted{Remarque}
Comment interpréter le foncteur
contravariant $\mathcal V^\circ-k\text{-mod gr}$ défini par $A$,
en composant $A$ et le foncteur $\omega^*:C^\circ\longrightarrow k\text{-mod gr}$
(associant à $(f_*,\unsure{f^*})$ l'application $f^*$)? En d'autres
termes, comment interpréter $A(f)^*:A(Y)-A(X)$,
si $f:Y\longrightarrow X$? On devine que c'est défini
par $\varphi(f)^*\in\operatorname{Hom}_{\mathcal V^A}(Y,X)=A(Y\times X)$, déduit
de $\varphi(f)\in\unsure{H|A}(X\times Y)$ par l'isom can. qu'on
devine. Vérification comme dessus.

On définit un anti-foncteur canonique
\[
\unsure{s}:\unsure{\mathcal V^{A\circ}}-\unsure{\mathcal V^{A\circ}},
\]
induisant l'identité sur les objets, par
\[
u\longmapsto{}^su:\operatorname{Hom}_{\mathcal V^A}(X,Y)=A(X\times Y)-\operatorname{Hom}_{\mathcal V^A}(Y,X)=A(Y\times X)
\]
\margin{vérifier que c'est \uline{\unsure{multiplicatif}!}}
\mspage{14}
